\documentclass[11pt,letterpaper]{article}
\usepackage[margin=1in]{geometry}
\usepackage{enumitem}
\usepackage{xcolor}
\usepackage{soul}
\usepackage{ulem}
\usepackage{titlesec}
\usepackage{parskip}
\usepackage{amssymb}

\definecolor{noteblue}{RGB}{70,130,180}
\definecolor{warnorange}{RGB}{210,105,30}
\definecolor{ideagreen}{RGB}{60,120,60}

\titleformat{\section}{\large\bfseries\color{noteblue}}{\thesection}{1em}{}
\titleformat{\subsection}{\normalsize\bfseries\color{warnorange}}{\thesubsection}{1em}{}

\setlength{\parindent}{0pt}

\begin{document}

\begin{center}
{\LARGE\textbf{THE PERSISTENCE}}\\[0.5em]
{\large Working Notes / Brain Dump / Style Bible}\\[0.3em]
{\small\textit{v0.3 --- not for distribution --- half-baked thoughts ahead}}
\end{center}

\vspace{1em}
\hrule
\vspace{1em}

\section*{CORE CONCEPT --- THE ELEVATOR PITCH}

Memory isn't storage, it's \textbf{transport}. You don't remember your wedding day, you \textit{go there}. Full sensory, full emotional, full phenomenological weight. The past isn't past---it's a place you can visit.

\textbf{BUT:} You still live forward. Time still moves. People still die. You just... never lose access to when they were alive.

The question isn't ``what if we could remember everything perfectly''---it's ``what happens to the \textit{present} when the past is always available?''

\vspace{1em}

\section*{TONE \& STYLE NOTES}

\begin{itemize}[leftmargin=*]
    \item \textbf{Quiet devastation} --- not dystopia-loud, not utopia-bright. The technology is \textit{normalized}. No one monologues about how amazing/terrible it is. It just \textit{is}.
    
    \item \textbf{Melancholy without melodrama} --- the sadness should creep, not announce itself. A woman laughing alone in a memorial hall. A man touching air where someone used to be.
    
    \item \textbf{Mundane future} --- fax machines still exist. Coffee still gets cold. Eggs still burn. The tech is revolutionary but the \textit{texture of life} remains recognizable.
    
    \item \textbf{Show the seams} --- don't over-explain. Let readers feel confused the way a time-traveler would feel confused. Trust them.
    
    \item \textbf{Verb tense as weapon} --- play with present/past tense to mirror the dual-presence experience. ``She was there and is there and will always be there'' etc.
\end{itemize}

\subsection*{What this is NOT:}

\begin{itemize}[leftmargin=*]
    \item Not \textit{Black Mirror} nihilism (technology bad! humans doomed!)
    \item Not techno-utopianism (we solved death! everything is fine!)
    \item Not a mystery/thriller about memory crimes (though those exist in-world)
    \item Not an action plot with blocking-weapons as MacGuffins
\end{itemize}

\vspace{1em}

\section*{WORLDBUILDING CAVEATS --- THINGS TO DECIDE}

\begin{enumerate}[leftmargin=*]
    \item \textbf{Is the tech biological, mechanical, or hybrid?}
    
    Current draft implies neural lacing + external storage. But is it:
    \begin{itemize}
        \item Implanted at birth? Opt-in surgery? Genetic modification?
        \item Universal or class-stratified? (draft says stratified but this has HUGE implications)
        \item Removable? What happens if the hardware fails?
    \end{itemize}
    
    \textcolor{ideagreen}{$\rightarrow$ LEAN INTO: the tech is so ubiquitous that its origins are fuzzy to most people. Like asking ``who invented WiFi'' --- technically answerable but culturally irrelevant.}
    
    \item \textbf{What about children?}
    
    If encoding starts at birth, do toddlers have accessible infant experiences? Is there a ``first memory'' anymore? What does childhood development look like when you can revisit your own tantrums?
    
    \textcolor{warnorange}{$\rightarrow$ DANGER ZONE: could get creepy fast. Maybe encoding requires neural maturity? First accessible memories around age 4-5?}
    
    \item \textbf{Legal/consent implications}
    
    \begin{itemize}
        \item Can your memories be subpoenaed?
        \item Can you delete experiences? Is deletion legal?
        \item What about witnesses to crimes---can they be \textit{forced} to revisit?
        \item Divorce proceedings where both parties have perfect recall of every fight???
    \end{itemize}
    
    \textcolor{ideagreen}{$\rightarrow$ RICH TERRITORY but don't let it dominate. This is literary fiction with SF elements, not legal thriller.}
    
    \item \textbf{The ``other people's memories'' problem}
    
    Draft mentions black market for accessing others' experiences. But HOW? If the encoding is neural, how do you transfer? 
    
    Options:
    \begin{itemize}
        \item Direct neural link (intimate, requires trust/coercion)
        \item Extracted/copied files (implies your experiences are \textit{data})
        \item Some kind of broadcast/reception (public vs private experiences)
    \end{itemize}
    
    \textcolor{warnorange}{$\rightarrow$ BE CAREFUL: this edges toward ``uploading consciousness'' territory which is a different story.}
\end{enumerate}

\vspace{1em}

\section*{THEMATIC THREADS TO PULL}

\subsection*{The Presence Problem}

If you can always go back, why be here?

\begin{itemize}[leftmargin=*]
    \item Addiction framing---but addiction to \textit{what}? Your own life?
    \item The present becomes a raw material factory for future visits
    \item ``Experience hoarding'' --- living moments \textit{for} the archive
    \item What does spontaneity mean when you know you'll revisit?
\end{itemize}

\textit{Possible scene: Someone on a first date who keeps thinking ``this will be a good one to revisit'' and thereby ruins it by not being present. The memory they create is of not-being-present. Forever.}

\subsection*{Grief Transformed (not eliminated)}

You never lose access to the dead. But:

\begin{itemize}[leftmargin=*]
    \item They stop \textit{generating} new moments
    \item The archive becomes finite, bounded, \textit{complete}
    \item You can visit them forever but they can't visit you
    \item Is it comfort or torture? (Answer: yes)
\end{itemize}

\textit{Possible scene: Someone realizes they've visited every moment with their dead spouse. There are no more surprises. The archive is a museum now, not a life.}

\subsection*{Identity \& Continuity}

If you can perfectly revisit yourself at 20, are you the same person?

\begin{itemize}[leftmargin=*]
    \item The ``younger self as stranger'' phenomenon
    \item Watching yourself make mistakes you can't undo
    \item Nostalgia for who you used to be vs. actually \textit{being} who you used to be
    \item Do people get ``stuck'' at certain ages? Refuse to update their self-concept?
\end{itemize}

\textcolor{ideagreen}{$\rightarrow$ This is where it gets genuinely philosophical. Don't answer, just explore.}

\subsection*{The Blocking Threat}

Loss of access = new form of death

\begin{itemize}[leftmargin=*]
    \item Not physical death, but experiential death
    \item Your past still exists but \textit{you can't get there}
    \item The terror isn't of forgetting (passive) but of being \textit{locked out} (active)
    \item Blocking as weapon, as punishment, as ``treatment''
\end{itemize}

\textit{Note: The university plot is about well-intentioned blocking (addiction treatment) weaponized by true believers. The horror is that BOTH SIDES think they're helping.}

\vspace{1em}

\section*{STRUCTURAL IDEAS}

\subsection*{Option A: Braided Timelines}

Multiple POV characters, each dealing with different aspects of the technology:

\begin{itemize}
    \item Mara \& Chen: Long-term relationship, navigating shared archives
    \item Kai: Someone who \textit{resists} visiting, tries to live only in present
    \item Ezra \& Yuki: Insiders who know what blocking really does
    \item Subject 7 (or equivalent): The human cost, shown not told
\end{itemize}

\textcolor{warnorange}{Risk: Could feel like interconnected short stories rather than novel. Needs strong throughline.}

\subsection*{Option B: Single POV, Expanding Scope}

One character (Mara?) whose personal story gradually reveals societal implications.

\begin{itemize}
    \item Starts intimate: her relationship, her visits, her choices
    \item Expands: she becomes aware of blocking, of the crisis, of the stakes
    \item Climax: personal and political collide
\end{itemize}

\textcolor{ideagreen}{More traditional but potentially more powerful. Readers experience the world through one lens.}

\subsection*{Option C: The Document Novel}

Excerpts, memos, transcripts, fragments. No single narrative.

\begin{itemize}
    \item University memos
    \item Therapy session transcripts
    \item Legal depositions
    \item Personal journal entries (what does a journal look like when you can just revisit?)
    \item News articles, academic papers, advertising copy
\end{itemize}

\textcolor{warnorange}{High risk/high reward. Could be brilliant or unreadable. Requires very strong voice consistency across formats.}

\vspace{1em}

\section*{LANGUAGE \& VOCABULARY}

\subsection*{Words that are OBSOLETE:}

\begin{itemize}
    \item Forgot / forgotten / forget
    \item Reminisce (replaced by ``visit'' or ``access'')
    \item Nostalgia (the emotion doesn't exist)
    \item ``Remember when...'' (now just ``when...'')
    \item Faded / hazy / dim (memories don't degrade)
    \item Deja vu (you KNOW if you've been there before)
\end{itemize}

\subsection*{NEW vocabulary:}

\begin{itemize}
    \item \textbf{Visiting} --- accessing a past experience
    \item \textbf{Weighted} / \textbf{weighting} --- how much attention is on present vs. past
    \item \textbf{Drift} / \textbf{drifting} --- unintentional visiting, losing present-focus
    \item \textbf{When} (as noun) --- a specific temporal location (``which when?'')
    \item \textbf{Archive} --- totality of one's accessible experiences
    \item \textbf{Temporal coordinates} --- specific timestamps for access
    \item \textbf{Blocking} --- suppression of revisitation access
    \item \textbf{Now-sickness} --- anxiety/depression related to present-focus
    \item \textbf{Presence} --- being fully weighted in current moment (rare, notable)
\end{itemize}

\subsection*{Phrases to AVOID:}

\begin{itemize}
    \item ``Back in my day'' --- still works but means something different
    \item ``If only I could go back'' --- you CAN, that's the point
    \item ``I'll never forget'' --- meaningless, nothing is forgotten
    \item ``Time heals all wounds'' --- does it though? when you can reopen them infinitely?
\end{itemize}

\vspace{1em}

\section*{QUESTIONS I DON'T HAVE ANSWERS TO YET}

\begin{enumerate}[leftmargin=*]
    \item What do \textbf{dreams} look like in this world? Draft says they're not captured, but people still dream---do dreams feel \textit{different} now? More ephemeral because they're the only ephemeral thing?
    
    \item What about \textbf{learning}? If you can revisit a lecture infinitely, does education change? Do skills still require practice or can you just revisit the moment you did it right?
    
    \item \textbf{Art and creativity} --- is there still a market for fiction when everyone has infinite real experiences to access? Or does fiction become MORE valuable as the only source of novelty?
    
    \item \textbf{Religion} --- how do concepts of afterlife, karma, redemption change when your life is already preserved eternally? Do people visit their own deaths? (Can they?)
    
    \item The \textbf{animals} question --- do pets have persistence? Can you access a dog's experience of you? (Probably not, but the desire would exist)
    
    \item \textbf{Sex and intimacy} --- obviously revisiting intimate moments is a thing, but what are the ethics? The etiquette? If you revisit sex with an ex, is that cheating? What if they're dead?
    
    \item \textbf{The singularity question} --- if experiences can be stored and accessed, are they uploadable? Transferable? Is this a step toward digital consciousness or fundamentally different?
\end{enumerate}

\vspace{1em}

\section*{STUFF I'M STEALING FROM / INFLUENCED BY}

\begin{itemize}[leftmargin=*]
    \item \textit{Eternal Sunshine of the Spotless Mind} --- memory as relationship, loss as identity
    \item Borges, ``Funes the Memorious'' --- the curse of perfect memory
    \item \textit{Station Eleven} --- quiet apocalypse, focus on human texture
    \item Ted Chiang generally --- rigorous SF with emotional core
    \item \textit{Never Let Me Go} --- normalized horror, characters who don't question their world
    \item Tarkovsky's \textit{Solaris} --- memory as haunting, presence of the dead
    \item Anne Carson's \textit{Nox} --- grief as archive, fragment as form
\end{itemize}

\vspace{1em}

\section*{IMMEDIATE NEXT STEPS}

\begin{enumerate}[leftmargin=*]
    \item[$\square$] Decide on structure (braided vs. single POV vs. documents)
    \item[$\square$] Write the blocking crisis from INSIDE --- what does it feel like?
    \item[$\square$] Develop Kai's ``presence resistance'' --- why does he refuse to visit?
    \item[$\square$] Figure out the ending --- is this tragedy? Hope? Ambiguity?
    \item[$\square$] Research: actual neuroscience of memory reconsolidation (for plausibility)
    \item[$\square$] Write a scene that's JUST two people visiting the same moment together
    \item[$\square$] The memorial hall needs more --- make it a full chapter?
\end{enumerate}

\vspace{2em}
\hrule
\vspace{1em}

\begin{center}
\textit{``The past is never dead. It's not even past.'' --- Faulkner}\\[0.5em]
\textit{But what if he was being literal?}
\end{center}

\end{document}
